\chapter{随机变量之和与大数定律}\label{ch:g12:sums}

为什么赌场最终总是获利，民意调查为何有效？因为
许多独立随机量的平均值波动越来越小。本章
证明这一点：期望的线性性、
独立变量的方差可加性、比内梅–切比雪夫不等式，
以及大数定律。

\section{随机变量之和}

设 $X, Y$ 为定义在同一有限样本空间
$\Omega$ 上的两个随机变量，则和 $X + Y$ 是随机变量
$\omega \mapsto X(\omega) + Y(\omega)$.

\begin{theorem}[期望的线性性]\label{thm:g12:sums:linearity}
\index{期望！线性性}
对 $X, Y$（均定义在 $\Omega$ 上）及 $a, b \in \R$，有：
\[
\E(X + Y) = \E(X) + \E(Y), \qquad \E(aX + b) = a\E(X) + b .
\]
\emph{此处无需假设独立性。}
\end{theorem}

\begin{proof}
将期望写成对结果的求和。由于
$\P(X = x) = \sum_{\omega : X(\omega) = x} \P(\{\omega\})$，合并各项后
得 $\E(X) = \sum_{\omega \in \Omega} \P(\{\omega\})\,X(\omega)$。于是
\[
\E(X + Y) = \sum_{\omega} \P(\{\omega\})\bigl(X(\omega) + Y(\omega)\bigr)
= \sum_{\omega} \P(\{\omega\})X(\omega)
+ \sum_{\omega} \P(\{\omega\})Y(\omega) = \E(X) + \E(Y). \qedhere
\]
\end{proof}

\begin{definition}[独立随机变量]\label{def:g12:sums:indep}
称 $X$ 与 $Y$ 为\emph{独立}\index{随机变量！独立}的，若
对所有取值 $x, y$ 有：
\[
\P(X = x \text{ 且 } Y = y) = \P(X = x)\,\P(Y = y).
\]
若干变量 $X_1, \dots, X_n$ 独立，是指
对每个子族中变量的每一组取值，乘积法则均成立。
\end{definition}

\begin{proposition}\label{prop:g12:sums:prodexp}
若 $X$ 与 $Y$ 独立，则 $\E(XY) = \E(X)\,\E(Y)$。
\end{proposition}

\begin{proof}
\[
\begin{aligned}
\E(XY) &= \sum_{x, y} xy\;\P(X = x \text{ 且 } Y = y)
= \sum_{x, y} xy\,\P(X=x)\P(Y=y) \\
&= \Bigl(\sum_x x\P(X=x)\Bigr)\Bigl(\sum_y y\P(Y=y)\Bigr). \qedhere
\end{aligned}
\]
\end{proof}

\begin{theorem}[和的方差]\label{thm:g12:sums:variance}
\index{方差！和的}
若 $X$ 与 $Y$ \emph{独立}，则
\[
\V(X + Y) = \V(X) + \V(Y).
\]
更一般地，对独立的 $X_1, \dots, X_n$，
$\V(X_1 + \dots + X_n) = \V(X_1) + \dots + \V(X_n)$。
\end{theorem}

\begin{proof}
利用柯尼希–惠更斯公式（\cref{prop:g12:randvar:konig}）及线性性：
\begin{align*}
\V(X + Y) &= \E\bigl((X+Y)^2\bigr) - \bigl(\E X + \E Y\bigr)^2\\
&= \E(X^2) + 2\E(XY) + \E(Y^2)
- \E(X)^2 - 2\E(X)\E(Y) - \E(Y)^2\\
&= \V(X) + \V(Y) + 2\bigl(\E(XY) - \E(X)\E(Y)\bigr),
\end{align*}
对独立变量，最后的括号项为零
（见\cref{prop:g12:sums:prodexp}）。一般情形由归纳法得到。
\end{proof}

\begin{example}[二项分布再探]\label{ex:g12:sums:binomial}
二项变量 $X \sim \mathcal B(n, p)$ 是
$X = X_1 + \dots + X_n$，为 $n$ 个独立伯努利变量之和。因此
从结构上可直接得出：
\[
\E(X) = np, \qquad \V(X) = np(1-p),
\]
无需任何计算即可重新得到\cref{thm:g12:randvar:binomial}。
\end{example}

\begin{proposition}[样本均值]\label{prop:g12:sums:mean}
设 $X_1, \dots, X_n$ 相互独立且与 $X$ 同分布
（期望为 $\mu$，方差为 $\sigma^2$），令
$M_n = \frac{X_1 + \dots + X_n}{n}$ 为它们的\emph{样本
均值}\index{样本均值}。则
\[
\E(M_n) = \mu, \qquad \V(M_n) = \frac{\sigma^2}{n}, \qquad
\sigma(M_n) = \frac{\sigma}{\sqrt n}.
\]
\end{proposition}

\begin{proof}
由线性性得 $\E(M_n) = \frac{n\mu}{n} = \mu$；由独立性得
$\V(X_1 + \dots + X_n) = n\sigma^2$，除以 $n$ 时方差
按 $\frac{1}{n^2}$ 缩放（见\cref{prop:g12:randvar:affine}）。
\end{proof}

$\frac{\sigma}{\sqrt n}$ 体现了基本的\emph{平方根律}：要
将平均值的波动减半，须使样本量变为原来的四倍。

\begin{omfigure}
\begin{tikzpicture}
\begin{axis}[omaxis,
  width=9.6cm, height=5.6cm,
  xlabel={$n$}, ylabel={$\sigma(M_n)/\sigma$},
  xmin=0, xmax=105, ymin=0, ymax=1.15,
  xtick={4,16,64}, ytick={1,0.5,0.25},
  yticklabels={$1$,$\tfrac12$,$\tfrac14$}]
  \addplot[omDef, thick, domain=1:100, samples=80] {1/sqrt(x)};
  \draw[black, dashed, thin] (axis cs:0,0.5) -- (axis cs:4,0.5)
    -- (axis cs:4,0);
  \draw[black, dashed, thin] (axis cs:0,0.25) -- (axis cs:16,0.25)
    -- (axis cs:16,0);
  \draw[black, dashed, thin] (axis cs:64,0.125) -- (axis cs:64,0);
  \fill (axis cs:1,1) circle (1.5pt);
  \fill (axis cs:4,0.5) circle (1.5pt);
  \fill (axis cs:16,0.25) circle (1.5pt);
  \fill (axis cs:64,0.125) circle (1.5pt);
\end{axis}
\end{tikzpicture}

{\small 平方根律：样本量每扩大为原来的四倍，
仅使均值的标准差减半。}
\end{omfigure}

\section{集中不等式}

\begin{theorem}[马尔可夫不等式]\label{thm:g12:sums:markov}
\index{马尔可夫不等式}
若 $X \geq 0$ 且 $a > 0$：
\[
\P(X \geq a) \leq \frac{\E(X)}{a}.
\]
\end{theorem}

\begin{proof}
在 $\E(X) = \sum_i p_i x_i$ 中（各项均非负），只保留
满足 $x_i \geq a$ 的项：每一项至少为 $a\,p_i$，于是
$\E(X) \geq a \sum_{x_i \geq a} p_i = a\,\P(X \geq a)$。
\end{proof}

\begin{theorem}[比内梅–切比雪夫不等式]\label{thm:g12:sums:chebyshev}
\index{比内梅–切比雪夫不等式}
对任意随机变量 $X$ 及任意 $\varepsilon > 0$：
\[
\P\bigl(\abs{X - \E(X)} \geq \varepsilon\bigr)
\leq \frac{\V(X)}{\varepsilon^2}.
\]
\end{theorem}

\begin{proof}
对非负随机变量 $Y = (X - \E(X))^2$ 应用
马尔可夫不等式，取 $a = \varepsilon^2$：
\[
\P\bigl(\abs{X - \E(X)} \geq \varepsilon\bigr)
= \P(Y \geq \varepsilon^2)
\leq \frac{\E(Y)}{\varepsilon^2} = \frac{\V(X)}{\varepsilon^2}. \qedhere
\]
\end{proof}

\begin{omfigure}
\begin{tikzpicture}
\begin{axis}[omaxis,
  width=9.6cm, height=5cm,
  ylabel={}, xmin=-3.4, xmax=3.4, ymin=0, ymax=1.25,
  xtick={-1,0,1},
  xticklabels={$\E(X)-\varepsilon$,$\E(X)$,$\E(X)+\varepsilon$},
  ytick=\empty]
  \addplot[name path=bell, omDef, thick, domain=-3.2:3.2]
    {exp(-x^2/2)};
  \path[name path=xaxis] (-3.2,0) -- (3.2,0);
  \addplot[omDef!20] fill between[of=bell and xaxis,
    soft clip={domain=-1:1}];
  \addplot[omThm!20] fill between[of=bell and xaxis,
    soft clip={domain=-3.2:-1}];
  \addplot[omThm!20] fill between[of=bell and xaxis,
    soft clip={domain=1:3.2}];
  \draw[black, dashed, thin] (axis cs:-1,0) -- (axis cs:-1,0.75);
  \draw[black, dashed, thin] (axis cs:1,0) -- (axis cs:1,0.75);
\end{axis}
\end{tikzpicture}

{\small 集中现象：比内梅–切比雪夫不等式给出
随机变量 $X$ 落入尾部（红色部分，与其期望的距离至少为 $\varepsilon$）的
概率的上界为 $\V(X)/\varepsilon^2$。}
\end{omfigure}

\begin{theorem}[大数定律]\label{thm:g12:sums:lln}
\index{大数定律}
设 $X_1, \dots, X_n$ 为独立同分布的
随机变量（期望为 $\mu$，方差为 $\sigma^2$），$M_n$ 为其样本均值。对
每个 $\varepsilon > 0$：
\[
\P\bigl(\abs{M_n - \mu} \geq \varepsilon\bigr)
\leq \frac{\sigma^2}{n\,\varepsilon^2}
\xrightarrow[n \to +\infty]{} 0 .
\]
样本均值集中于期望附近。
\end{theorem}

\begin{proof}
对 $M_n$ 应用比内梅–切比雪夫不等式，其期望为 $\mu$，
方差为 $\frac{\sigma^2}{n}$（见\cref{prop:g12:sums:mean}）。
\end{proof}

\begin{remark}
本定理是概率论与统计学之间的桥梁：
在多次独立重复中，事件的\emph{频率}趋近于其
\emph{概率}（取 $X_i$ 为该事件的示性变量，从而
$\mu = p$）。它为用模拟估计概率
（\emph{蒙特卡罗方法}）以及用民意调查估计总体比例 ---
定量分析见\cref{ch:g12:contdist}。
\end{remark}

\begin{method}[运用比内梅–切比雪夫不等式]\label{met:g12:sums:chebyshev}
为保证 $\P(\abs{M_n - \mu} \geq \varepsilon) \leq \alpha$，
只需 $n \geq \frac{\sigma^2}{\alpha\,\varepsilon^2}$ 即可。该界较粗（实际波动通常小得多），
但十分普适：
它除有限方差外不要求分布的任何其他信息。
\end{method}

\section{练习}

\begin{exercise}[$\star$]\label{exo:g12:sums:1}
掷两枚均匀骰子，设 $S$ 为点数之和。
利用线性性（而非 $S$ 的分布！），计算 $\E(S)$；
再利用独立性计算 $\V(S)$，已知单枚均匀骰子的方差为 $\frac{35}{12}$。
\end{exercise}

\begin{exercise}[$\star$]\label{exo:g12:sums:2}
设 $X \sim \mathcal B(100,\ 0.5)$。用马尔可夫不等式估计
$\P(X \geq 75)$ 的上界，再用比内梅–切比雪夫不等式估计
$\P(\abs{X - 50} \geq 25)$ 的上界，并加以比较。
\end{exercise}

\begin{exercise}[$\star$]\label{exo:g12:sums:3}
掷一枚均匀硬币 $n$ 次，记 $F_n$ 为正面出现的频率。
问：$n$ 至少要多大，才能由比内梅–切比雪夫界推出
$\P\bigl(\abs{F_n - 0.5} \geq 0.05\bigr) \leq 0.05$ 成立？
\end{exercise}

\begin{exercise}[$\star\star$]\label{exo:g12:sums:4}
一位投资者将其资金平均分配给 $n$ 项相互独立的资产，
每项资产的期望收益率为 $\mu = 5\%$，标准差为
$\sigma = 20\%$。求组合收益率 $M_n$ 的期望与标准差，
以及使标准差低于 $4\%$ 所需的资产数目，
并指出这一例子体现了
何种金融原理？
\end{exercise}

\begin{exercise}[$\star\star$]\label{exo:g12:sums:5}
设 $X$ 与 $Y$ 相互独立，且都在 $\{1, 2, 3\}$ 上服从均匀分布。
\begin{enumerate}
  \item 写出 $S = X + Y$ 的分布，
        并由此直接计算 $\E(S)$、$\V(S)$。
  \item 利用期望的线性性与方差的可加性，重新求出这两个值。
\end{enumerate}
\end{exercise}

\begin{exercise}[$\star\star$]\label{exo:g12:sums:6}
说明在无独立性时，方差的可加性可能不成立：
计算 $\V(X + Y)$（取 $Y = X$），并与 $\V(X) + \V(Y)$ 比较。
哪些随机变量 $X$ 使得 $\V(2X) = 2\V(X)$ 成立？
\end{exercise}

\begin{exercise}[$\star\star$]\label{exo:g12:sums:7}
怀疑一枚骰子不均匀。将其掷 $1200$ 次，
出现六点 $260$ 次（$f = 0.2167$，而非 $\frac16 \approx 0.1667$）。
在骰子均匀的假设下，估计
$\P\bigl(\abs{F_n - \frac16} \geq 0.05\bigr)$ 的上界（用比内梅–切比雪夫不等式），
并讨论该均匀性假设是否合理。
\end{exercise}

\begin{exercise}[$\star\star\star$]\label{exo:g12:sums:8}
\emph{（关于硬币的一个更好的不等式。）}
设 $X \sim \mathcal B(n,\ p)$ 且 $F_n = \frac Xn$。
\begin{enumerate}
  \item 证明：$p(1-p) \leq \frac14$ 对 $p \in \intcc{0}{1}$ 成立。
  \item 由此推出与分布无关的上界
        $\P\bigl(\abs{F_n - p} \geq \varepsilon\bigr) \leq
        \frac{1}{4n\varepsilon^2}$。
  \item 至少需要调查多少人，才能使观测频率
        与真实比例之差不超过 $3$ 个百分点的概率至少为
        $95\%$，用此界估计？（实际民意调查使用更精细的估计，
        但数量级是正确的。）
\end{enumerate}
\end{exercise}

\section{问题：庄家总是赢家}

\begin{problem}[{周末问题 --- 大数定律解释赌场、民意调查与保险，
并顺带驳斥了
赌徒谬误}]
\label{pb:g12:sums:1}
一名轮盘赌玩家转 $100$ 次后，几乎有一半的时候
处于\emph{领先}。转动一百万次的赌场，则
以任何法庭都无可质疑的确定性保持领先。游戏相同，优势同样微小 ---
差别在于 $\sqrt n$，这正是本章的主题
（见\cref{prop:g12:sums:mean}、
\cref{thm:g12:sums:chebyshev}、\cref{thm:g12:sums:lln}）。
本题核算赌场账目、确定选举民意调查的规模、
为分散化投资定价 --- 并驳斥赌博史上代价最高的
谬误。

\medskip
\noindent\textbf{第一部分 --- 熟练运算。}
\begin{enumerate}
  \item 两枚独立的骰子：计算 $V(X + Y)$
        （见\cref{thm:g12:sums:variance}）。
  \item 掷 $100$ 枚独立的骰子并取平均值：
        求该样本均值的期望与标准差
        （见\cref{prop:g12:sums:mean}；对单枚骰子，
        $\sigma = \sqrt{35/12} \approx 1.71$）。
  \item 一个非负随机变量的期望为 $2$。
        由马尔可夫不等式
        （见\cref{thm:g12:sums:markov}），
        对 $\P(X \geq 10)$ 能给出什么结论？
  \item 估计 $\P\left(\abs{\bar X_{100} - 3.5} \geq
        0.5\right)$ 的上界（取第 2 小题中骰子的平均值），
        用比内梅–切比雪夫不等式。
  \item 期望总是可加的，而方差只在
        独立时才可加。举出不独立的 $X, Y$（提示：
        $Y = -X$）使得
        $V(X + Y) \neq V(X) + V(Y)$。
\end{enumerate}

\medskip
\noindent\textbf{第 II 部分 --- 庄家的账本。}
欧式轮盘，押注 $1$ 欧元于红色：赢得 $+1$ 的
概率为 $\frac{18}{37}$，输掉 $-1$ 的概率为
$\frac{19}{37}$。
\begin{enumerate}[resume]
  \item 计算一次押注
        收益的期望与标准差。
  \item 某赌徒进行 $100$ 次相互独立的押注，设 $G$ 为
        总收益。计算 $\E(G)$ 与 $\sigma(G)$，再求
        比值 $\frac{\abs{\E(G)}}{\sigma(G)}$。问：
        标准差四分之一的漂移量，对赌徒
        今晚保持领先的机会意味着什么？
  \item 赌场共接受 $1\,000\,000$ 次押注。计算
        赌场总盈利的期望与标准差，
        以及同样的比值。解释数 $27$ 的含义。
  \item 用切比雪夫不等式证明：估计赌场在
        百万次押注中\emph{亏钱}的概率的上界。
  \item 押注策略：输后加倍，领先时停止\,\dots\
        利用期望的线性性
        （\cref{thm:g12:sums:linearity}），对（可能随机的）
        押注序列，说明为何\emph{每一种}
        策略在负优势博弈中的期望收益都为负
        ——一个能获胜的系统会违背什么？
  \item 叙述大数定律
        （\cref{thm:g12:sums:lln}）对红色的频率
        作出了何种承诺，以及它\emph{不}承诺什么，
        即连续十次红色之后的下一次转动。指出
        该谬误的名称。
  \item 最微妙之处：大数定律靠的是\emph{稀释}，
        而不是补偿。若某晚红色比期望多出 $100$，
        这部分盈余永远不会
        “被偿还”——而应改为计算
        \emph{比例} $\frac{100}{n}$ 在 $n = 10^6$ 时的变化，
        并用一句话改写赌徒谬误的错误。
\end{enumerate}

\medskip
\noindent\textbf{第 III 部分 --- 民意调查。}
\begin{enumerate}[resume]
  \item 由\cref{exo:g12:sums:8}，至少需要调查多少人，
        才能使观测频率与真值之差在 $3$ 个百分点之内
        的概率至少为 $95\,\%$？
        用与分布无关的切比雪夫界估计？
  \item 实际民调机构通常用约 $1\,100$ 人，以获得
        $\pm 3$ 个百分点、置信度 $95\,\%$ 的误差范围，其公式为
        $n \approx \frac{1.96^2 \times p(1-p)}
        {\varepsilon^2}$，其中 $p(1-p) \leq \frac14$。代入计算
        它的值，并解释与第 13 题的差距（切比雪夫
        不等式对波动的形状
        一无所知）。
  \item 要将误差范围减半，样本量须怎样增长？
        从误差 $\pm 3$ 个百分点、样本量 $1\,100$ 出发，
        何种样本量能给出 $\pm 1.5$ 个百分点？
  \item 一个反直觉的经典结论：所需样本量从未用到
        总体的规模——$1\,100$ 人即可用于一座城市
        或一整个大洲。指出模型中
        缺少总体规模的那一处，并给出民调机构常用的
        “厨房类比”。
\end{enumerate}

\medskip
\noindent\textbf{第 IV 部分 --- 分散投资。}
\begin{enumerate}[resume]
  \item \cref{pb:g12:randvar:1}中的保险公司
        期望收益为 $30n$ 欧元（$n$ 为保单份数），
        理赔的标准差约为 $315\sqrt n$。
        对于哪些 $n$，期望利润最终超过
        理赔的一个标准差？
        $\sqrt n$ 对保险公司起了什么作用？
  \item 一只股票的年收益率波动
        幅度为 $\sigma = 20\,\%$。将资金平均分配于 $25$
        只\emph{相互独立}的此类股票，计算该组合的
        $\sigma$。金融界把分散投资
        称为唯一的“免费午餐”——这顿午餐究竟是什么？
  \item 现设这 $25$ 只股票完全相关（它们
        同涨同跌）：该组合的 $\sigma$ 为多少？
        比较两种极端情形，并指出每一种分散化论证
        暗中依赖的是哪一个假设——
        以及 2008 年该假设在全系统范围失效时发生了什么。
  \item 终曲——$\sqrt n$ 交响曲：和的波动
        幅度为 $\sqrt n$，而平均数则以 $\frac{1}{\sqrt n}$ 的
        速度趋于稳定；请将这一主题贯穿本题的四个行业
        即赌场、民调、保险与投资组合，
        指出两个长期存在的“滥用者”（赌徒谬误和
        虚假的独立性），并给出
        前瞻：波动的\emph{形状}
        ——钟形曲线——是下一章连续情形的
        主角，其完整定理将为大学分册
        压轴。
\end{enumerate}
\end{problem}
